\begin{figure}[Diagrama de componentes del SCDEE]%
              {fig:scdee-components}%
              {Diagrama de componentes del backend del SCDEE.}
  \image{}{}{mermaid/scdee_components}
\end{figure}

El sistema se descompone en siete componentes principales
(\cref{fig:scdee-components}):

\begin{enumerate}

  \item \textbf{Servicio \acs{api} (Django + \acs{drf}).}
  Aplicación web que expone los recursos del sistema bajo el prefijo
  \texttt{/api/v1/}, implementa el \textit{middleware} de contexto
  multi-tenant, la autenticación basada en \dfnpl{plugin} y la
  autorización contextual por \texttt{SubjectMembership}; delega toda
  tarea de larga duración a los \textit{workers} asíncronos. La
  versión~1.0 utiliza el servidor de desarrollo de Django
  (\texttt{runserver}); en un despliegue de producción este se sustituye
  por un servidor \acs{wsgi} dedicado como Gunicorn o uWSGI, sustitución
  transparente al código de la aplicación.

  \item \textbf{PostgreSQL.}
  Base de datos relacional principal, que aloja toda la información
  estructurada del dominio (organizaciones, usuarios, asignaturas,
  exámenes, instancias, calificaciones, anotaciones, notificaciones, log
  de auditoría) con integridad referencial estricta y transacciones
  \acs{acid} (\cref{rnf:persistence}).

  \item \textbf{Redis.}
  Base de datos en memoria que actúa simultáneamente como \textit{broker}
  de Celery, \textit{result backend}, caché de la configuración por
  organización y almacén de la lista negra de \textit{refresh tokens}
  revocados.

  \item \textbf{MinIO.}
  Servicio de almacenamiento de objetos compatible con la \ac{api}
  \ac{s3} que aloja los binarios de tamaño variable: \acsp{pdf} en blanco
  de los modelos, \acsp{pdf} instrumentados con códigos \acs{qr}, páginas
  escaneadas individuales (\texttt{ExamPage}), composiciones bajo demanda
  y \textit{payloads} de anotaciones (\cref{rnf:object-storage}).

  \item \textbf{\textit{Workers} de Celery.}
  Procesos asíncronos suscritos a tres colas
  (\texttt{default}, \texttt{recognition}, \texttt{notifications}) que
  ejecutan respectivamente las tareas de mantenimiento, reconocimiento
  de páginas y envío de correo. Las tres colas se sirven actualmente
  desde el mismo proceso para simplificar el despliegue local, pero la
  arquitectura admite separarlas en \textit{workers} dedicados sin
  cambios de código en producción.

  \item \textbf{Celery beat.}
  Planificador de tareas periódicas con \textit{backend} en base de
  datos (mediante el paquete \texttt{django-celery-beat}), responsable
  de disparar el borrado automático de exámenes archivados
  (\cref{rf:config:auto-purge}), la apertura y el cierre de las ventanas
  de revisión (\cref{rf:review:auto-open,rf:review:auto-close}) y otras
  tareas programadas.

  \item \textbf{\textit{Watcher} de ingesta.}
  Comando de gestión de Django (\texttt{python manage.py sftp\_watcher})
  que se ejecuta como proceso aparte y monitoriza una o varias carpetas
  locales en busca de nuevas páginas escaneadas. Al detectar un fichero
  válido, lo reclama atómicamente, codifica su contenido y encola la
  tarea de ingesta en Celery. La organización a la que pertenece cada
  página se extrae del código \acs{qr} durante el reconocimiento, no
  de la estructura del directorio. En el despliegue local del SCDEE las
  carpetas monitorizadas son volúmenes Docker; en un despliegue real
  con cabina de escaneo serían carpetas \acs{sftp} expuestas por el
  escáner.

\end{enumerate}

A estos siete componentes se añade un servicio \acs{smtp} externo, utilizado
por los \textit{workers} para el envío asíncrono de correo (en desarrollo
se sustituye por una instancia local de Mailpit que captura los mensajes
sin entregarlos), y un conjunto opcional de servicios de monitorización
(pgAdmin, Redis Commander y Flower) que el operador puede activar bajo
demanda mediante un perfil aparte de Docker Compose. El \dfn{frontend},
responsabilidad del \acs{tfg} paralelo, se sitúa fuera del backend y se
comunica con él exclusivamente a través de la \acs{api} \acs{rest}; esta
separación permite que ambos extremos evolucionen y se desplieguen de forma
independiente.

Dentro del servicio \acs{api}, el código se organiza siguiendo una
arquitectura por capas con una capa de servicios explícita
(\textit{Service Layer}, descrita por Fowler \cite{fowler2002}): los
modelos representan la persistencia, los \textit{serializers} la validación
y transformación de entrada y salida, las vistas el enrutado y la
autorización, y los servicios la lógica de negocio reutilizable y
desacoplada de \acs{http}. La \cref{SEC:DES-LAYERS} desarrolla en detalle
las responsabilidades de cada capa.
